\section{A highly transitive action with finite-type stabilizers}
\label{sec:stage}

We strengthen the embedding problem to a statement about an action and
all its finite-point stabilizers.

\begin{theorem}\label{thm:action}
For every finitely generated group $G$ with decidable word problem,
there is a group $\Lambda$ containing $G$ and acting faithfully on a
countably infinite set $S$ such that:
\begin{enumerate}[label=\textup{(\roman*)}]
 \item the action is transitive on ordered tuples of distinct points
 of every finite length;
 \item $\Lambda$ and every pointwise stabilizer of a finite tuple
 are of type $\Finfty$.
\end{enumerate}
\end{theorem}

The action will be obtained by iterating an embedding of a group and its
set of columns.  Finite permutations become available after one
transition, giving high transitivity in the limit.  We then identify each
finite-point stabilizer as an ascending torus and verify the two
hypotheses of Theorem~\ref{thm:torus}: commuting splittings and a
factorization through a finitely presented group.

\subsection{Ring data and compressed diagonal matrices}

Apply Theorem~\ref{thm:ring-data} to the given finitely generated group
$G$ with decidable word problem.  Throughout this section, retain its
ring $B$, injective additive multiplicative map $\phi$, idempotent
$p=\phi(1)$, and centralizing binary pairs $s_i,d_i\in pBp$:
\begin{equation}\label{eq:stage-binary}
 d_i s_j=\delta_{ij}p,
 \qquad s_0d_0+s_1d_1=p.
\end{equation}
In particular, $\phi$ need not preserve the identity.  Also retain the
nonzero corner $eBe$, its embedded groups $G$ and $W$, and its elements
$S_y,T_y$ for $y\in Y$.  Write
\[
 Q=\E_7(B),\qquad X=B^7,
\]
and let
$\iota:Q\hookrightarrow W$ and $\iota_X:X\hookrightarrow Y$
be the equivariant embeddings of Theorem~\ref{thm:ring-data}.
For $x\in X$ we abbreviate
$S_{\iota_X(x)}$ and $T_{\iota_X(x)}$ to $S_x$ and $T_x$.

The first use of the binary corner is to turn a unit of $B$ into an
elementary matrix supported in one coordinate.  This will let us encode
both the action of $W$ and finite permutations of its indexed prefixes
inside $Q$.

For a unit $v\in B^\times$, denote by $D_j(v)$ the diagonal matrix
with entry $v$ in position $j$ and identity entries elsewhere; its size
will always be specified by the surrounding group.  Write
$e_{ij}(b)=I+b\mathbf e_{ij}$ for an elementary matrix, where
$\mathbf e_{ij}$ is a matrix unit.

\begin{lemma}\label{lem:diagonal}
For every $r\ge2$ and $v\in B^\times$,
\[
 D_1\bigl(1+\phi(v-1)\bigr)\in\E_r(B).
\]
\end{lemma}

\begin{proof}
Let $\operatorname{GE}_r(B)$ be the subgroup of $\GL_r(B)$ generated
by elementary matrices and invertible diagonal matrices.  Conjugation
by a diagonal matrix preserves the elementary subgroup, so
$\E_r(B)$ is normal in $\operatorname{GE}_r(B)$.  The quotient is
abelian, even though $B$ need not be commutative.  Indeed, in a
two-coordinate block,
\[
 w(v):=e_{12}(v)e_{21}(-v^{-1})e_{12}(v)
   =\begin{pmatrix}0&v\\-v^{-1}&0\end{pmatrix},
 \qquad
 w(v)w(-1)=\begin{pmatrix}v&0\\0&v^{-1}\end{pmatrix}.
\]
Thus $D_i(v)$ and $D_j(v)$ have the same class modulo $\E_r(B)$.
Classes of one-place diagonals commute because they can be represented
in distinct positions.  These classes generate the quotient.

Set
\[
 E_i=s_i d_i\quad(i=0,1),
 \qquad
 E_{00}=s_0s_0d_0d_0,
 \qquad
 E_{01}=s_0s_1d_1d_0.
\]
The binary relations give orthogonal decompositions
\begin{equation}\label{eq:stage-four-corners}
 p=E_0+E_1,
 \qquad E_0=E_{00}+E_{01}.
\end{equation}
All these idempotents commute with $\phi(B)$.

Each of $E_0,E_{00},E_{01}$ is orthogonal to and equivalent to $E_1$,
with implementing elements that commute with $\phi(B)$.  Here are
explicit implementing elements.  Put
\[
 s_{00}=s_0s_0,\quad d_{00}=d_0d_0,
 \qquad
 s_{01}=s_0s_1,\quad d_{01}=d_1d_0.
\]
For $v\in\{0,00,01\}$, define
$a_v=s_1d_v$ and $b_v=s_vd_1$.  Then
\[
 a_vb_v=E_1,\qquad b_va_v=E_v,
 \qquad
 \tau_v=1-E_v-E_1+a_v+b_v
\]
satisfies $\tau_v^2=1$ and $\tau_vE_v\tau_v^{-1}=E_1$.
It also commutes with $\phi(B)$.

Fix a unit $u\in B^\times$.  For any of these idempotents $E$,
the element $1+\phi(u-1)E$ is a unit, with inverse
$1+\phi(u^{-1}-1)E$.  Conjugating by $\tau_v$ shows that the
one-place diagonals belonging to $E_v$ and $E_1$ have equal classes
in the abelian quotient $\operatorname{GE}_r(B)/\E_r(B)$.
Let $c$ be the common class for $E_0,E_1,E_{00},E_{01}$.
The second orthogonal split in
Equation~\eqref{eq:stage-four-corners} gives $c=c^2$, hence $c=1$.
The first split then shows that
$D_1(1+\phi(u-1)p)$ is elementary.  Since
$\phi(u-1)p=\phi(u-1)$, this is the asserted matrix.
\end{proof}

\subsection{The transition of acting pairs}

The corner embedding of $W$ defines an injective homomorphism
\[
 \rho:W\hookrightarrow B^\times,
 \qquad \rho(w)=1-e+w.
\]
Indeed, its inverse-unit formula is
$\rho(w)^{-1}=1-e+w^{-1}$, and multiplication is the multiplication
in $eBe$ together with the identity on $1-e$.
Write $\phi^2=\phi\circ\phi$ and define
\begin{align}
 \mu(w)&=D_1\bigl(1+\phi^2(\rho(w)-1)\bigr)
        \in Q,\label{eq:stage-mu}\\
 \beta(y)&=\bigl(\phi^2(S_y),0,\ldots,0\bigr)^{\mathsf T}
        \in X.\label{eq:stage-beta}
\end{align}
Membership of $\mu(w)$ in $Q$ follows from
Lemma~\ref{lem:diagonal}, applied twice.  Expanding a product shows
that $v\mapsto1+\phi(v-1)$ is a homomorphism on units, and it is
injective because $\phi$ is injective.  Together with the injectivity
of $\rho$, this proves that $\mu$ is an injective homomorphism.

The prefix relations show that $\beta$ is injective.  If $y\ne z$
and $S_y=S_z$, multiplication by $T_y$ would give $e=0$;
injectivity of $\phi^2$ preserves the resulting distinction.
They also give equivariance, since
\begin{align*}
 \bigl(1+\phi^2(w-e)\bigr)\phi^2(S_y)
  &=\phi^2(S_y)+\phi^2\bigl((w-e)S_y\bigr)\\
  &=\phi^2(wS_y)
   =\phi^2(S_{wy}).
\end{align*}
We thus obtain an injective transition of acting pairs
\begin{equation}\label{eq:stage-transition}
 \delta=\mu\iota:Q\hookrightarrow Q,
 \qquad
 \sigma=\beta\iota_X:X\hookrightarrow X.
\end{equation}
The same compressed-diagonal construction, applied to the corner
copy of $G$, gives an embedding $G\hookrightarrow Q$.

\begin{lemma}\label{lem:stage-permutations}
Every permutation of $\sigma(X)$ with finite support is induced
on that set by an element of $Q$.
\end{lemma}

\begin{proof}
Let $F\subseteq X$ be finite, and let $\eta$ be a permutation of
$F$.  Put
\[
 v=1-\sum_{x\in F}S_xT_x
      +\sum_{x\in F}S_{\eta(x)}T_x.
\]
The elements $S_xT_x$ are pairwise orthogonal idempotents.  The
relations $T_xS_y=\delta_{xy}e$ show directly that $v$ is a unit,
with inverse obtained by replacing $\eta$ by $\eta^{-1}$.
Moreover,
\[
 vS_x=S_{\eta(x)}\quad(x\in F),
 \qquad vS_x=S_x\quad(x\in X\setminus F).
\]
The matrix $D_1(1+\phi^2(v-1))$ lies in $Q$ by
Lemma~\ref{lem:diagonal}, and acts on $\sigma(X)$ as required.
\end{proof}

\subsection{The limit group and its action}

Let
\[
 N=\varinjlim(Q,\delta),
 \qquad S=\varinjlim(X,\sigma).
\]
Write $q_i$ for the element of $N$ represented by $q\in Q$ at
stage $i$, and $[x,i]$ for a point of $S$ represented by $x\in X$.
Thus
\[
 q_i=\delta(q)_{i+1},
 \qquad [x,i]=[\sigma(x),i+1].
\]
The maps in both direct systems are injective.  Hence their stages
may be identified with increasing subgroups $Q_i\le N$ and
increasing subsets $X_i\subseteq S$.
Equivariance in Equation~\eqref{eq:stage-transition} defines an
action of $N$ on $S$: for $k\ge i,j$,
\begin{equation}\label{eq:stage-limit-action}
 q_i[x,j]=
 \bigl[\delta^{k-i}(q)\sigma^{k-j}(x),k\bigr].
\end{equation}
Passing to a larger common stage leaves this expression unchanged.

The shift $\tau(q_i)=q_{i+1}$ is an automorphism of $N$, with
inverse $\tau^{-1}(q_i)=\delta(q)_i$.  Similarly the bijection
\[
 t[x,i]=[x,i+1]
\]
of $S$ has inverse $[x,i]\mapsto[\sigma(x),i]$.
Equation~\eqref{eq:stage-limit-action} implies that this bijection
conjugates the action of $n\in N$ to that of $\tau(n)$.
It therefore extends the action to
\[
 \Lambda=N\rtimes_{\tau}\langle t\rangle.
\]
Identifying $Q$ with $Q_0$, its defining relation is
\begin{equation}\label{eq:stage-torus}
 t^{-1}qt=\delta(q)\qquad(q\in Q).
\end{equation}
In particular, $\Lambda$ is the ascending torus of $\delta$.
For completeness, the inverse map from the direct-limit description
to the group with this presentation sends $q_i$ to $t^iqt^{-i}$.
Equation~\eqref{eq:stage-torus} makes this well defined and gives
the two inverse homomorphisms.

\begin{lemma}\label{lem:stage-faithful-transitive}
The set $S$ is countably infinite.  The action of $\Lambda$ on $S$
is faithful and transitive on ordered tuples of distinct points
of every finite length.
\end{lemma}

\begin{proof}
The finitely generated ring $B$ is countable, so $X$ and $S$ are
countable.  The ring is infinite: iterating the binary relations
in Equation~\eqref{eq:stage-binary} produces, for every $m$, a
family of $2^m$ pairwise orthogonal nonzero idempotents in $pBp$.
They are nonzero because the corresponding deletion-prefix products
are $p\ne0$.  Thus $X$, and consequently $S$, is infinite.

The action of $N$ is faithful.  Indeed, a nonidentity element of
$N$ is represented by a nonidentity matrix in some stage $Q_i$.
Such a matrix moves one of the standard basis columns in that stage,
and injectivity of the point maps preserves the distinction in $S$.

Every element of $N$ preserves $X_j$ setwise for all sufficiently
large $j$: its action at any such stage is by an invertible matrix.
In contrast, $tX_j=X_{j+1}$, and every stage inclusion is strict.
To see strictness, note that $\sigma(X)$ has zero coordinates
outside position 1, whereas the column with identity in position 2
does not.  Thus, for $k\ne0$ and sufficiently large $j$,
$t^kX_j=X_{j+k}\ne X_j$.  No nonzero power of $t$ acts as an
element of $N$.  If an element $nt^k\in\Lambda$ acts trivially,
this observation forces $k=0$, and faithfulness for $N$ then forces
$n=1$.

Given two ordered tuples of distinct points of the same finite
length, place them in a common stage $X_i$.  Their images in the
next stage lie in $\sigma(X)$.  The prescribed bijection between
their entries extends to a finitely supported permutation of
$\sigma(X)$.
Lemma~\ref{lem:stage-permutations} supplies the required element
of $Q_{i+1}$.  This proves the asserted transitivity.
\end{proof}

\subsection{Pointwise stabilizers as ascending tori}

Fix a nonnegative integer $m$.  By
Lemma~\ref{lem:stage-faithful-transitive}, it suffices to prove the
finiteness assertion for one ordered tuple $z=(z_1,\ldots,z_m)$
of distinct points with every $z_j\in\sigma(X)\subseteq X$.
When $m=0$ the tuple is empty and its stabilizer is all of
$\Lambda$.  Since both $z$ and $\sigma(z)$ have entries in
$\sigma(X)$, Lemma~\ref{lem:stage-permutations} gives $a\in Q$
such that
\begin{equation}\label{eq:stage-align}
 a\sigma(z)=z.
\end{equation}
For the empty tuple, take $a=1$.
Put
\[
 U=Q_{(z)},\qquad t'=ta^{-1}.
\]
The subscript denotes pointwise stabilization.  In our convention
for the limit action,
\[
 t'[z,0]=[a^{-1}z,1]=[\sigma(z),1]=[z,0],
\]
so $t'$ fixes the tuple.  Furthermore,
\begin{equation}\label{eq:stage-stabilizer-map}
 (t')^{-1}qt'=a\delta(q)a^{-1},
 \qquad
 f:U\longrightarrow U,
 \quad f(u)=a\delta(u)a^{-1}.
\end{equation}
The codomain in this formula follows either from the fact that $t'$
fixes the tuple or directly from equivariance and
Equation~\eqref{eq:stage-align}.  The map $f$ is injective.

\begin{lemma}\label{lem:stage-stabilizer-torus}
The pointwise stabilizer $\Lambda_{([z,0])}$ is the ascending torus
of $U$ under $f$.
\end{lemma}

\begin{proof}
The height homomorphism $\Lambda\to\Z$ sends $Q$ to zero and
$t'$ to one.  Its kernel is $N$.  For every $i\ge0$ there is
$b_i\in Q$ such that $(t')^i=t^i b_i$: this follows inductively
from $qt=t\delta(q)$.  Hence
\[
 (t')^iQ(t')^{-i}=t^iQt^{-i}=Q_i,
 \qquad N=\bigcup_{i\ge0}(t')^iQ(t')^{-i}.
\]
Since $t'$ fixes the tuple,
\[
 N\cap\Lambda_{([z,0])}
   =\bigcup_{i\ge0}(t')^iU(t')^{-i}.
\]
An element of the full stabilizer of height $k$, multiplied by
$(t')^{-k}$, belongs to this height-zero subgroup.  Thus the full
stabilizer is its semidirect product with $\langle t'\rangle$.
The conjugation formula in
Equation~\eqref{eq:stage-stabilizer-map} identifies this group
with the ascending torus of $f$.
\end{proof}

The two diagram systems required by Theorem~\ref{thm:torus}
also restrict to $U$.  Let
$j_{\ell,E}:Q\to W$ be the maps of
Theorem~\ref{thm:ring-data}, for $\ell=2,3$ and
$E\in\mathcal P_\ell$, with common full-identity map $\iota$.
Define
\begin{equation}\label{eq:stage-stabilizer-diagrams}
 h_{\ell,E}:U\longrightarrow U,
 \qquad
 h_{\ell,E}(u)=a\mu(j_{\ell,E}(u))a^{-1}.
\end{equation}
These maps do have their entire images in $U$.  If $u$ fixes $z_j$,
the fixed-point property of the enlargement says that
$j_{\ell,E}(u)$ fixes $\iota_X(z_j)$.  Equivariance of $\beta$
then says that $\mu(j_{\ell,E}(u))$ fixes $\sigma(z_j)$;
conjugation by $a$ gives the claim.  The maps are homomorphisms,
$h_{\ell,I}=f$, and they preserve each orthogonal-sum identity
and each required elementwise commutation relation, since $\mu$
and conjugation by $a$ are homomorphisms.

We have now identified the stabilizer as an ascending torus and supplied
its commuting splittings.  The remaining hypothesis of
Theorem~\ref{thm:torus} is a factorization of $f$ through a finitely
presented group.  We will use $\St_6(B)$ as that group.  A six-coordinate
elementary block inside $Q=\E_7(B)$ fixes every vector supported in the
first coordinate, so its entire image fixes $\sigma(z)$.  Two steps are
needed: lift the compressed units homomorphically to $\St_6(B)$, then
move their nonidentity part from coordinate 1 into that block without
moving the marked columns.

\subsection{A homomorphic Steinberg lift}

For a unital ring $R$, the Steinberg group $\St_6(R)$ has generators
$x_{ij}(a)$ for $i\ne j$ and $a\in R$, with the relations
\begin{align*}
 x_{ij}(a)x_{ij}(b)&=x_{ij}(a+b),\\
 [x_{ij}(a),x_{jk}(b)]&=x_{ik}(ab)
       &&\text{if }i,j,k\text{ are distinct},\\
 [x_{ij}(a),x_{kl}(b)]&=1
       &&\text{if }i\ne l\text{ and }j\ne k.
\end{align*}
We use the convention $[u,v]=uvu^{-1}v^{-1}$.
The map sending $x_{ij}(a)$ to $e_{ij}(a)$ is a surjection
$\pi:\St_6(R)\twoheadrightarrow\E_6(R)$.

We need the following precise form of the kernel-annihilation theorem
from \citet[Theorem~4.1]
{FoundationsCompanion}.

\begin{theorem}[Kernel annihilation]\label{thm:stage-kernel}
Let $R$ be a unital associative ring and let $\psi:R\to R$ be an
additive multiplicative map.  Put $P=\psi(1)$.  Suppose that
$U_0,U_1,V_0,V_1\in PRP$ commute with $\psi(R)$ and satisfy
\[
 V_iU_j=\delta_{ij}P,
 \qquad U_0V_0+U_1V_1=P.
\]
Then the homomorphism
\[
 \theta:\St_6(R)\longrightarrow\St_6(R),
 \qquad x_{ij}(a)\longmapsto x_{ij}(\psi(a)),
\]
annihilates the kernel of $\pi$.
\end{theorem}

The theorem permits nonunital $\psi$ and noncommutative $R$.
For such a coefficient map the matrix of the image of a root word
with matrix $g$ is
\begin{equation}\label{eq:stage-nonunital-matrix}
 I+\psi^{(6)}(g-I),
\end{equation}
where $\psi^{(6)}$ means entrywise application.  This follows first
on elementary matrices and then on products from
$gh-I=(g-I)+(h-I)+(g-I)(h-I)$.

\begin{lemma}\label{lem:steinberg-lift}
There is a homomorphism
$L:B^\times\to\St_6(B)$ such that
\begin{equation}\label{eq:stage-lift}
 \pi L(v)=D_1\bigl(1+\phi^2(v-1)\bigr).
\end{equation}
Moreover, $\St_6(B)$ is finitely presented.
\end{lemma}

\begin{proof}
Apply Theorem~\ref{thm:stage-kernel} with
\[
 R=B,\quad \psi=\phi,\quad P=p,
 \quad U_i=s_i,\quad V_i=d_i.
\]
Its hypotheses are precisely the support, commutation, and binary
relations in Theorem~\ref{thm:ring-data}.
Consequently its coefficient map factors as a homomorphism
\[
 \lambda:\E_6(B)\longrightarrow\St_6(B),
 \qquad \lambda(\pi(s))=\theta(s).
\]
The definition is independent of the choice of $s$ because $\theta$
kills $\ker\pi$.  By
Equation~\eqref{eq:stage-nonunital-matrix},
\[
 \pi\lambda(g)=I+\phi^{(6)}(g-I).
\]

The homomorphism on units used to define $\mu$, together with
Lemma~\ref{lem:diagonal}, supplies a homomorphism
\[
 A:B^\times\longrightarrow\E_6(B),
 \qquad A(v)=D_1\bigl(1+\phi(v-1)\bigr).
\]
Take $L=\lambda A$.  The displayed formula for $\pi\lambda$ gives
Equation~\eqref{eq:stage-lift}.

Finally, Lemma~\ref{lem:encoding-finite-presentation} proves that $B$
is finitely presented as a unital associative ring: to its finite
$\F_2$-algebra presentation one adds $2\cdot1=0$ over $\Z$.
Theorem~3 of
\citet{KrsticMcCool1999} states that $\St_n(R)$ is finitely
presented for every finitely presented unital associative ring $R$
and every $n\ge4$.  It allows noncommutative rings and applies with
$R=B$ and $n=6$.
\end{proof}

\subsection{Factoring the stabilizer endomorphism}

The tuple $\sigma(z)$ is supported in coordinate 1, as is the
nonidentity part of $\delta(U)$.  An ordinary exchange of coordinates
1 and 2 would move the tuple.  Instead we exchange only the idempotent
part on which those nonidentity increments are supported; this part
annihilates the marked columns.  The resulting factorization sends the
whole intermediate Steinberg group into $U$.

\begin{lemma}\label{lem:stage-stabilizer-factor}
The endomorphism $f$ in
Equation~\eqref{eq:stage-stabilizer-map} factors through
$\St_6(B)$ by homomorphisms
\[
 U\xrightarrow{\alpha}\St_6(B)\xrightarrow{b}U.
\]
\end{lemma}

\begin{proof}
For the fixed tuple $z$, put
\begin{equation}\label{eq:stage-d}
 q_j=\phi^2(S_{z_j}T_{z_j}),
 \qquad d=1-\sum_{j=1}^m q_j,
 \qquad c=1-d.
\end{equation}
The prefix relations make the $q_j$ pairwise orthogonal
idempotents, so $d$ is an idempotent.  It depends only on $z$.

For $u\in U$, write $w=\iota(u)$ and
\[
 r=\phi^2(\rho(w)-1)=\phi^2(w-e),
 \qquad v'=1+r.
\]
Then $\delta(u)=D_1(v')$.  The element $w$ fixes each index
$\iota_X(z_j)$, so both prefix equations give
\[
 (w-e)S_{z_j}=0,
 \qquad T_{z_j}(w-e)=0.
\]
Consequently $rq_j=q_jr=0$ and
\begin{equation}\label{eq:stage-support}
 dr=rd=r.
\end{equation}
Also, using $T_{z_i}S_{z_j}=\delta_{ij}e$,
\begin{equation}\label{eq:stage-fixed-columns}
 d\phi^2(S_{z_j})
  =\phi^2(S_{z_j})
       -\sum_{i=1}^m\phi^2(S_{z_i}T_{z_i}S_{z_j})=0.
\end{equation}
These identities use only additivity and multiplicativity of $\phi$;
no identity-preservation assumption is needed.

In coordinates 1 and 2 of $B^7$, let
\[
 E=e_{12}(d)e_{21}(-d)e_{12}(d)
   =\begin{pmatrix}c&d\\-d&c\end{pmatrix},
 \qquad
 E^{-1}=\begin{pmatrix}c&-d\\d&c\end{pmatrix},
\]
with identity entries on all other coordinates.  This is an element
of $Q$.  Equation~\eqref{eq:stage-fixed-columns} shows that both
$E$ and $E^{-1}$ fix every column $\sigma(z_j)$.  Direct
multiplication, using Equation~\eqref{eq:stage-support}, gives
\begin{equation}\label{eq:stage-exchange}
 E D_1(v')E^{-1}=D_2(v').
\end{equation}
In particular, the signed exchange leaves no sign or inversion on
the moved unit.

Let $L$ be the homomorphism of Lemma~\ref{lem:steinberg-lift},
and define
\begin{equation}\label{eq:stage-alpha}
 \alpha(u)=L(\rho(\iota(u))).
\end{equation}
By Equation~\eqref{eq:stage-lift}, its matrix projection, placed
in coordinates 2 through 7, is $D_2(v')$.  In particular, the
argument of $L$ here is the original unit $\rho(w)$; the two
applications of $\phi$ are already included in the defining
property of $L$.

For every $s\in\St_6(B)$, define
\begin{equation}\label{eq:stage-b}
 b(s)=aE^{-1}\diag(1,\pi(s))Ea^{-1}.
\end{equation}
This is a homomorphism to $Q$, because embedding the elementary
six-coordinate block in coordinates 2 through 7 gives an elementary
seven-coordinate matrix.  Its entire image lies in $U$: the
matrix $\diag(1,\pi(s))$ fixes every first-coordinate vector
$\sigma(z_j)$, the matrices $E^{\pm1}$ do so as well, and
$a\sigma(z)=z$.  Finally,
\[
 b(\alpha(u))
   =aE^{-1}D_2(v')Ea^{-1}
   =aD_1(v')a^{-1}
   =f(u).
\]
This proves the factorization.  For $m=0$, the sums in
Equation~\eqref{eq:stage-d} are empty and $d=1$; all the same
formulas apply.

The codomains can be seen in the commutative diagram
\[
\begin{tikzcd}[column sep=large,row sep=large]
 U \arrow[r,"\alpha"] \arrow[d,"f"']
   & \St_6(B) \arrow[d,"{\diag(1,\pi(-))}"] \\
 U
   & Q_{(\sigma(z))} \arrow[l,"c_{aE^{-1}}"']
\end{tikzcd}
\]
where $c_g$ denotes conjugation by $g$. The lower right group is
the entire pointwise stabilizer of $\sigma(z)$ in $Q$.
The right arrow lands there because it leaves the first coordinate
unchanged; the lower arrow lands in $U$ because $E$ fixes
$\sigma(z)$ and $a\sigma(z)=z$.
\end{proof}

\begin{proof}[Proof of Theorem~\ref{thm:action}]
The embedding $G\hookrightarrow Q\hookrightarrow\Lambda$ was
constructed above, and
Lemma~\ref{lem:stage-faithful-transitive} proves the assertions
about the action.  For a tuple $z$ chosen as above,
Lemma~\ref{lem:stage-stabilizer-torus} identifies its pointwise
stabilizer with the ascending torus of the injective endomorphism
$f:U\to U$.  The diagrams in
Equation~\eqref{eq:stage-stabilizer-diagrams} satisfy the two
orthogonal-splitting hypotheses of Theorem~\ref{thm:torus}.
Lemma~\ref{lem:stage-stabilizer-factor}, together with finite
presentation in Lemma~\ref{lem:steinberg-lift}, supplies its other
hypothesis.  Thus this stabilizer has type $\Finfty$.
High transitivity conjugates every tuple of distinct points to one
of this form.  Repeated entries do not change a pointwise
stabilizer, and the empty tuple gives $\Lambda$ itself.  This
proves all the claims.
\end{proof}
